Zu jedem Punkt
\mathl{P \in M}{} wählen wir eine offene Kartenumgebung
\mathl{P \in U_P}{} mit einer Karte
\maabbdisp {\alpha_P} {U_P } {V_P
} {}
mit
\mathl{V_P \subseteq \R^d}{.} Dabei können wir annehmen, indem wir zu einer Ballumgebung von
\mathl{\alpha_P(P) \in V_P}{} und dessen Urbild übergehen, dass die $V_p$ offene Bälle sind, deren Radius maximal
\mathl{{ \frac{   1 }{  3 } }}{} ist. Die

\mathbed {U_P} {}
 {P\in M} {}
 {} {} {} {,}
überdecken die Mannigfaltigkeit. Wegen der Kompaktheit von $M$ gibt es endlich viele Punkte
\mathl{P_1 , \ldots , P_n}{} derart, dass auch

\mathbed {U_i=U_{P_i}} {}
 {i=1 , \ldots , n} {}
 {} {} {} {,}
die Mannigfaltigkeit überdecken. Wir platzieren die offenen Bälle
\mathl{V_i=V_{P_i}}{} in
\zusatzklammer {\anfuehrung{einem neuen}{}} {} {}
$\R^d$, und zwar mit den Mittelpunkten

\mathdisp {(1,0 , \ldots , 0), (2,0 , \ldots , 0) , \ldots , (n,0 , \ldots , 0)} { . }

Wegen der Radiusbedingung sind diese Bälle zueinander disjunkt. Wir betrachten die beschränkte offene Menge
\mathl{U= \bigcup_{i=1}^n V_i}{.} Die Kartenabbildungen $\alpha_i$ liefern stetige Abbildungen

\mathdisp {\varphi_i :V_i \stackrel{ \left(  \alpha_i  \right)^{-1} }{\longrightarrow} U_i \longrightarrow M} { . }

Wegen der Disjunktheit ergibt sich daraus durch
\mathl{\varphi(z)= \varphi_i(z)}{,} falls
\mathl{z \in V_i}{} ist, eine stetige Abbildung
\maabbdisp {\varphi} { U } {M
} {.}
Wegen der Überdeckungseigenschaft ist diese surjektiv.

 [[Kategorie:Latexseite]]
